\documentclass[a4paper]{article}
\usepackage[T1]{fontenc}
\usepackage[utf8]{inputenc}
\usepackage[brazil,shorthands=off]{babel}
\usepackage{newtxtext,newtxmath}
\usepackage[varqu]{zi4}
\usepackage[a4paper,margin=4mm]{geometry}
\usepackage{multicol,enumitem,tikz,fancyvrb,tcolorbox}
\usepackage[table]{xcolor}
\pagestyle{empty}
\setlength{\parindent}{0pt}
\setlength{\parskip}{0.5pt}
\setlength{\columnsep}{2.4mm}
\setlength{\columnseprule}{0.2pt}

% cores (escuras o bastante p/ sair legível em P&B)
\definecolor{hd}{RGB}{22,52,110}
\definecolor{tm}{RGB}{20,80,150}
\definecolor{ok}{RGB}{0,115,40}
\definecolor{wr}{RGB}{175,30,30}
\definecolor{nd}{RGB}{150,70,0}
\definecolor{tr}{RGB}{110,110,110}
\definecolor{cs}{RGB}{200,100,0}
\definecolor{es}{RGB}{20,100,200}
\definecolor{rg}{RGB}{130,40,150}
\definecolor{shadecolor}{RGB}{238,241,247}
\definecolor{bx}{RGB}{255,243,230}
\definecolor{oRL}{RGB}{0,125,125}
\definecolor{oWL}{RGB}{205,65,0}
\definecolor{oUL}{RGB}{40,90,210}
\definecolor{oC}{RGB}{0,120,0}
\definecolor{oA}{RGB}{190,0,0}
\definecolor{oLog}{RGB}{140,70,20}
\definecolor{kw}{RGB}{25,60,150}
\definecolor{fn}{RGB}{130,40,150}

\setlist{nosep,leftmargin=2.6mm,labelsep=0.7mm}
\setlist[itemize]{label={\color{hd}\raisebox{0.15ex}{\scalebox{0.7}{$\blacktriangleright$}}}}
\newenvironment{dl}{\begin{description}[leftmargin=3.2mm,labelindent=0pt,itemsep=0.7pt,topsep=0.6pt,font=\color{tm}\bfseries,style=standard]}{\end{description}}
\newcommand{\h}[1]{\par\vspace{2.2pt}{\setlength{\fboxsep}{0.9pt}\colorbox{hd}{\parbox{\dimexpr\linewidth-1.8pt}{\color{white}\bfseries\ #1}}}\par\vspace{0.8pt}}
\newcommand{\sub}[1]{\par\vspace{1.4pt}{\color{hd}\bfseries\underline{#1}}\par}
\newcommand{\ans}[1]{\leavevmode{\color{ok}\bfseries #1}}
\newcommand{\at}[1]{{\color{wr}\bfseries #1}}
\newcommand{\CS}{{\color{cs}\bfseries conservativo}}
\newcommand{\ES}{{\color{es}\bfseries estrito}}
\newcommand{\RG}{{\color{rg}\bfseries rigoroso}}
\newcommand{\go}{$\rightarrow$}
\newcommand{\Ra}{$\Rightarrow$}
\newcommand{\la}{$\leftarrow$}
\newcommand{\q}[1]{{\ttfamily #1}}
\newcommand{\sel}{$\sigma$}
\newcommand{\pj}{$\pi$}
\newcommand{\jn}{$\bowtie$}
\newcommand{\un}{$\cup$}
\newcommand{\N}[1]{\textcolor{nd}{\textbf{#1}}}
\newcommand{\Tr}[1]{\textcolor{tr}{#1}}
\newcommand{\K}[1]{\textcolor{kw}{\textbf{#1}}}
\newcommand{\F}[1]{\textcolor{fn}{\textbf{#1}}}
\newcommand{\oRL}[1]{\textcolor{oRL}{#1}}
\newcommand{\oWL}[1]{\textcolor{oWL}{\textbf{#1}}}
\newcommand{\oUL}[1]{\textcolor{oUL}{#1}}
\newcommand{\oW}[1]{\textbf{#1}}
\newcommand{\oC}[1]{\textcolor{oC}{\textbf{#1}}}
\newcommand{\oA}[1]{\textcolor{oA}{\textbf{#1}}}
\newcommand{\oLog}[1]{\textcolor{oLog}{\textbf{#1}}}
\newcommand{\aviso}[1]{\par\vspace{1pt}{\setlength{\fboxsep}{1.2pt}\colorbox{bx}{\parbox{\dimexpr\linewidth-2.4pt}{#1}}}\par\vspace{1pt}}
\fvset{fontsize=\fontsize{6.1}{6.75}\selectfont,commandchars=\\\{\}}
\newtcolorbox{sql}{colback=shadecolor,colframe=shadecolor,boxrule=0pt,arc=0pt,left=1.2pt,right=0.5pt,top=0.8pt,bottom=0.8pt,boxsep=0pt,before skip=1.5pt,after skip=1.5pt}
% grafico de trancas x tempo
\newcommand{\gr}[3]{\begin{tikzpicture}[x=0.95mm,y=0.85mm,baseline=0]
\draw[->,very thin](0,0)--(17,0);\draw[->,very thin](0,0)--(0,6.5);
\draw[densely dashed,very thin](13,0)--(13,6);\node[font=\fontsize{4}{4}\selectfont,anchor=south] at (13,5.6){C};
\draw[thick,#3]#2;\node[font=\fontsize{4.8}{4.8}\selectfont\bfseries,anchor=north,text=#3] at (8,0){#1};\end{tikzpicture}}

\begin{document}
\fontsize{7.35}{8.15}\selectfont\raggedright
\begin{multicols*}{3}
%==================== FRENTE: TRANSACOES ====================
{\setlength{\fboxsep}{1.2pt}\colorbox{bx}{\parbox{\dimexpr\linewidth-2.4pt}{\bfseries P2 (simulado): 4 questões múltipla escolha, 80 min: consulta distrib. rede social (3), 2PC (2 + 2), variante 2PL (3). A P1 teve Parte 2 c/ upload de PDF, a P2 pode ter tb. Provas dif. sorteadas.}}}

\h{ESCALONAMENTOS (base P1, cai junto)}
\begin{dl}
\item[conflito:] Ts dif., msm item, $\ge$1 W (RW, WR, WW)
\item[grafo preced.:] Ti\go Tj se op de Ti conflita c/ op \emph{posterior} de Tj. ciclo \Ra\ ñ seriável (0 seriais). s/ ciclo \Ra\ seriável, seriais = ordens topológicas
\item[recuperável:] T q leu item escrito por T' só confirma DPS de T'.\\ viola \q{R1(X)W1(X)R2(X)W2(X)C2}
\item[ACA:] (livre de rollback em cascata) ninguém LÊ item escrito por T' ñ confirmada (s/ leitura suja).\\ viola \q{R1(X)W1(X)R2(X)W2(X)}
\item[estrito:] ninguém LÊ NEM ESCREVE item escrito por T' ñ confirmada.\\ viola \q{R2(X)R1(X)W1(X)W2(X)}
\item[hierarquia:] estrito $\subset$ ACA $\subset$ recup. Leu sujo e os 2 abortam, ou a de quem leu confirma antes \go\ recuperável OK
\end{dl}
\sub{E3} base \q{R2(Y) R1(Y) R1(X) W1(X) R2(X) W2(Y)} + final:
\begin{itemize}
\item \q{..A1 A2} rec \ans{V} \quad \q{..C1 A2} rec \ans{V} \quad \q{..A1 C2} rec \ans{F}
\item \q{..C2 C1} rec \ans{F}, ACA \ans{F} \quad \q{..C2 A1} rec \ans{F}
\item \q{R2(Y) R1(Y) R1(X) W1(X) C1 R2(X) W2(Y) C2}: rec/ACA/estrito/seriável \ans{V} (\ans{T1\go T2})
\item \q{R2(Y) R1(Y) R1(X) R2(X) W2(Y) C2 W1(X) C1}: Y T1\go T2, X T2\go T1 \go\ ciclo, \ans{ñ seriável, nenhum serial}
\end{itemize}
\sub{Ex. 3.2} serial / seriável / nº seriais / rec / ACA / estrito
\begin{itemize}
\item \q{r1(X) r2(Z) r3(X) r1(Z) r2(Y) r3(Y) w1(X) c1 w2(Z) w3(Y) w2(Y) c3 c2}\\ \ans{N / N} (ciclo T2$\leftrightarrow$T3) \ans{/ 0 / S / S / N} (w2(Y) antes de c3)
\item \q{r1(X) r2(Z) r1(Z) r3(X) r3(Y) w1(X) c1 w3(Y) c3 r2(Y) w2(Z) w2(Y) c2}\\ \ans{N / S / 1} (T3\go T1\go T2) \ans{/ S / S / S}
\item \q{r1(X) r2(Z) r1(Z) r3(X) r3(Y) w1(X) w3(Y) r2(Y) w2(Z) w2(Y) c1 c2 c3}\\ \ans{N / S / 1} (T3\go T1\go T2) \ans{/ N} (T2 leu de T3, confirmou antes) \ans{/ N / N}
\end{itemize}

\h{TRANCAS (locks)}
\begin{dl}
\item[binária:] L(X)/UL(X), 0 livre 1 trancado. L e UL = região crítica (senão 2 T pegam a msm)
\item[compartilhada:] RL leitura (várias T), WL escrita (exclusiva), UL. RL espera se WL de outra; WL espera se WL ou RL de outra
\item[compat.:] RL+RL ok, resto ñ. upgrade RL\go WL pode (pode suspender)
\item[diretrizes:] WL antes de W; RL ou WL antes de R; soltar tudo antes de terminar (C ou A)
\item[obs:] só tranca ñ serializa: \q{RL1 R1 UL1 RL2 R2 UL2 WL1 W1 ...} continua ñ seriável
\end{dl}

\h{2PL (two-phase locking)}
\begin{dl}
\item[regra:] nenhuma T libera tranca antes de alocar TODAS q precisa. expansão (só aloca) \go\ \emph{lock point} \go\ contração (só solta)
\item[garante:] seriável (ordem dos lock points). Ñ evita deadlock (só \CS) nem cascata (só \ES/\RG)
\end{dl}
\begin{dl}
\item[\CS:] aloca TODAS, no estado + forte (vai ler e escrever X \go\ já pede WL(X)), ANTES da 1ª op R/W. Falta uma \go\ solta as q pegou e espera. s/ deadlock, pode starvation
\item[\ES:] WL só solta DPS do C/A (RL pode soltar antes). escal. estrito, recup. simples
\item[\RG:] TODA tranca (RL e WL) só solta dps do C. rigoroso \Ra\ estrito
\item[resumo:] \CS\ = QUANDO ALOCA (início); \ES/\RG\ = QUANDO SOLTA (fim). combinam
\end{dl}
\vspace{2pt}
\gr{simples}{(0,0)--(2,2)--(4,2)--(6,4)--(8,5)--(10,3.5)--(11.5,2)--(12.5,1)--(12.5,0)}{black!75}\hfill
\gr{conservativo}{(0,0)--(0.6,5)--(8,5)--(9.5,3)--(11,1.5)--(12,0)}{cs}\hfill
\gr{estrito}{(0,0)--(2,2)--(4,3)--(6,5)--(8,5)--(9,3.5)--(13,3.5)--(13,0)}{es}\\[2pt]
\gr{rigoroso}{(0,0)--(2,2)--(4,3)--(6,5)--(13,5)--(13,0)}{rg}\hfill
\gr{conserv.+rig.}{(0,0)--(0.6,5)--(13,5)--(13,0)}{cs}\hfill
\parbox[b]{17mm}{\fontsize{5}{5.6}\selectfont y = nº trancas de 1 T\\ x = tempo, C = commit\\ estrito: o q fica até C são as WL}
\vspace{1pt}
\sub{RECEITA: qual variante?} (p/ CADA T, em CADA nodo)
\begin{enumerate}[label=\textbf{\color{hd}\arabic*.},leftmargin=3mm]
\item RL/WL da T DPS de algum UL dela? \go\ \at{ñ é 2PL}
\item todos RL/WL antes do 1º R/W dela? \go\ \CS\\ (RL, opera e só dps WL no msm item = upgrade no meio = ñ conservativo)
\item todo UL de item c/ WL vem dps do C? \go\ \ES
\item TODO UL dps do C? \go\ \RG
\item 1 T q quebra já derruba a variante p/ o escal. todo
\end{enumerate}
\aviso{\at{READY/PREPARE ñ é commit!} UL antes do READY = UL antes do C \go\ ñ estrito. Log termina s/ UL da WL = ainda segurando \go\ estrito.}
\sub{Ex. slide 1:} \q{R1(X) R2(Y) W1(X) C1 W2(Y) C2}
\begin{dl}
\item[simples:] \q{RL1(X) R1(X) RL2(Y) R2(Y) WL1(X) W1(X) UL1(X) C1 WL2(Y) W2(Y) UL2(Y) C2}
\item[\CS:] \q{WL1(X) R1(X) WL2(Y) R2(Y) W1(X) UL1(X) C1 W2(Y) UL2(Y) C2}
\item[\ES/\RG:] \q{RL1(X) R1(X) RL2(Y) R2(Y) WL1(X) W1(X) C1 UL1(X) WL2(Y) W2(Y) C2 UL2(Y)}
\item[conserv.+estr.:] \q{WL1(X) R1(X) WL2(Y) R2(Y) W1(X) C1 UL1(X) W2(Y) C2 UL2(Y)}
\end{dl}
\sub{E3 Q15-18:} base \q{R2(Z) R1(Z) R1(Y) W1(Y) C1 R2(Y) W2(Z) C2}
\begin{dl}
\item[simples:] \q{RL2(Z) R2(Z) RL1(Z) R1(Z) RL1(Y) R1(Y) WL1(Y) UL1(Z) W1(Y) UL1(Y) C1 RL2(Y) R2(Y) WL2(Z) UL2(Y) W2(Z) UL2(Z) C2}
\item[\CS\ (ñ estr.):] \q{WL2(Z) RL2(Y) R2(Z) RL1(Z) R2(Y) UL2(Y) W2(Z) UL2(Z) C2 RL1(Z) WL1(Y) R1(Z) UL1(Z) R1(Y) W1(Y) UL1(Y) C1}
\item[conserv.+rig.:] \q{WL2(Z) RL2(Y) R2(Z) RL1(Z) R2(Y) W2(Z) C2 UL2(Y) UL2(Z) RL1(Z) WL1(Y) R1(Z) R1(Y) W1(Y) C1 UL1(Z) UL1(Y)}
\item[\ES\ (ñ cons.):] \q{RL2(Z) R2(Z) RL1(Z) R1(Z) RL1(Y) R1(Y) WL1(Y) UL1(Z) W1(Y) C1 UL1(Y) RL2(Y) R2(Y) WL2(Z) UL2(Y) W2(Z) C2 UL2(Z)}
\end{dl}
\sub{Ex. slide 3:} \q{R2(Y) R1(Y) R1(X) W1(X) C1 R2(X) W2(Y) C2}
\begin{dl}
\item[simples:] \q{RL2(Y) R2(Y) RL1(Y) R1(Y) RL1(X) R1(X) WL1(X) UL1(Y) W1(X) UL1(X) C1 RL2(X) R2(X) WL2(Y) UL2(X) W2(Y) UL2(Y) C2}
\item[\CS:] \q{WL2(Y) RL2(X) R2(Y) RL1(Y)}[susp.] \q{R2(X) UL2(X) W2(Y) UL2(Y) C2 RL1(Y) WL1(X) R1(Y) UL1(Y) R1(X) W1(X) UL1(X) C1}
\item[\ES:] \q{RL2(Y) R2(Y) RL1(Y) R1(Y) RL1(X) R1(X) WL1(X) UL1(Y) W1(X) C1 UL1(X) RL2(X) R2(X) WL2(Y) UL2(X) W2(Y) C2 UL2(Y)}
\item[\RG:] \q{RL2(Y) R2(Y) RL1(Y) R1(Y) RL1(X) R1(X) WL1(X) W1(X) C1 UL1(Y) UL1(X) RL2(X) R2(X) WL2(Y) W2(Y) C2 UL2(X) UL2(Y)}
\item[conserv.+estr.:] \q{WL2(Y) RL2(X) R2(Y) RL1(Y) R2(X) UL2(X) W2(Y) C2 UL2(Y) RL1(Y) WL1(X) R1(Y) UL1(Y) R1(X) W1(X) C1 UL1(X)}
\item[conserv.+rig.:] \q{WL2(Y) RL2(X) R2(Y) RL1(Y) R2(X) W2(Y) C2 UL2(X) UL2(Y) RL1(Y) WL1(X) R1(Y) R1(X) W1(X) C1 UL1(Y) UL1(X)}
\end{dl}

\h{DEADLOCK / STARVATION}
\begin{dl}
\item[deadlock:] ciclo no grafo de espera (Ti\go Tj: Ti espera tranca de Tj).\\ ex \q{RL1(X) R1(X) RL2(X) R2(X) WL1(X) WL2(X)}
\item[prevenir:] 2PL \CS; trancar itens sempre na msm ordem; por TS (menor = + velha):
\begin{itemize}
\item \textbf{wait-die}: velha pede \go\ espera; nova pede \go\ morre (reinicia c/ msm TS)
\item \textbf{wound-wait}: velha pede \go\ aborta a nova; nova pede \go\ espera
\item \textbf{no-waiting}: aborta direto. \textbf{cautious}: só espera se quem tem ñ está esperando
\end{itemize}
\item[detectar:] grafo de espera, achou ciclo \go\ aborta vítima
\item[timeout:] esperou demais \go\ aborta
\item[starvation:] T nunca roda (sempre vítima) \go\ prioridade / FIFO / reinicia c/ msm TS
\end{dl}

\h{OUTROS MÉTODOS DE CONCORRÊNCIA}
\begin{dl}
\item[timestamp (TO):] s/ trancas \Ra\ s/ deadlock. ts por T (ordem de início); Read\_TS(X), Write\_TS(X) = maior ts q leu/escreveu X. Aborta T se:
\begin{itemize}
\item R(X) c/ ts(T) $<$ Write\_TS(X)
\item W(X) c/ ts(T) $<$ Write\_TS(X) ou ts(T) $<$ Read\_TS(X)
\end{itemize}
rollback (até cascata) e starvation. TO estrito: espera quem escreveu X confirmar. Thomas: W obsoleto só é ignorado
\item[multiversão:] guarda versões antigas; leitura pega a versão q mantém seriável; rollback = descarta versões novas; + memória; extremo = BD temporal; c/ TS ou 2PL (snapshot: Postgres, Oracle)
\item[otimista:] executa livre (cópia local) \go\ validação \go\ escrita; falhou na validação \go\ aborta e reexecuta (slide BDD: desfaz a + nova). Bom c/ pouca interferência
\item[granularidade:] campo $<$ registro $<$ tabela $<$ BD; fina = + concorrência + overhead; intencionais IS, IX, SIX
\end{dl}

\h{ISOLAMENTO SQL-92}
{\rowcolors{2}{shadecolor}{white}\begin{tabular}{@{}l@{\hspace{2.5mm}}c@{\hspace{2mm}}c@{\hspace{2mm}}c@{}}
\textbf{nível} & \textbf{dirty} & \textbf{nonrep.} & \textbf{phantom}\\
READ UNCOMMITTED & sim & sim & sim\\
READ COMMITTED & -- & sim & sim\\
REPEATABLE READ & -- & -- & sim\\
SERIALIZABLE & -- & -- & --\\
\end{tabular}}
\begin{dl}
\item[dirty:] lê valor de T ñ confirmada
\item[nonrep.:] relê a linha e mudou (UPDATE de outra)
\item[phantom:] refaz consulta e aparece/some linha (INSERT/DELETE)
\item[default:] Postgres/Oracle READ COMMITTED, MySQL REPEATABLE READ
\item[SQL:] \q{SET TRANSACTION ISOLATION LEVEL ..; BEGIN; SAVEPOINT s; ROLLBACK TO s; COMMIT}. Cmd s/ BEGIN = transação por default
\end{dl}

\h{RECUPERAÇÃO DE FALHAS}
\begin{dl}
\item[falhas:] de T (erro, deadlock, user) / de sistema (perde cache, ñ disco) / de disco: backup+log = demorado; \ans{espelho sempre atualizado} = rápido (necessário e suficiente)
\item[steal:] grava pág. suja antes do C \go\ precisa UNDO
\item[no-force:] ñ força no C \go\ precisa REDO. SGBD real = steal+no-force = UNDO/REDO
\item[WAL:] reg. de log gravado ANTES do dado da cache ir p/ disco; log da T no disco antes do C
\item[adiada:] dado no disco só dps do C; log de W só valor \textbf{NOVO}; \ans{NO-UNDO/REDO}
\item[imediata:] dado pode ir antes do C; log: id T, id item, \textbf{ANTIGO e NOVO}; \ans{UNDO/REDO} (c/ force: UNDO/NO-REDO)
\item[shadow:] grava em \textbf{OUTRO lugar} do disco, troca diretório no C; \ans{NO-UNDO/NO-REDO}
\item[checkpoint:] força págs sujas p/ disco + marca; T c/ C antes do ckpt ñ precisa REDO
\item[receita:] lê log até último ckpt; separa confirmadas (têm C) e ativas; UNDO ativas de trás p/ frente (valor antigo); REDO confirmadas p/ frente (valor novo). UNDO e REDO \textbf{idempotentes}
\end{dl}
\sub{L2}
\begin{itemize}
\item desempenho mantendo recup. \go\ \ans{logs de várias T confirmadas em batch}
\item WAL \go\ \ans{log de cada op gravado antes dos dados em cache}
\item cascata qdo T \ans{lê item escrito por T ñ confirmada} (= def. de ACA; tirei 0,5 marcando só essa, talvez penalidade de tentativa)
\item adiada guarda \ans{novo}; imediata \ans{id T + id item + antigo + novo}; falha c/ adiada \go\ \ans{REDO}
\item refazer tudo desde backup \go\ \ans{demorado}; sombreamento \go\ \ans{local diferente}; disco rápido \go\ \ans{espelho}; UNDO/REDO \go\ \ans{idempotentes}
\end{itemize}

\h{SQL RÁPIDO}
\begin{sql}
\begin{Verbatim}
SELECT [DISTINCT] cols, COUNT(..)
FROM t1 [INNER JOIN t2 ON t1.a = t2.a]
WHERE cond          \Tr{-- filtra LINHAS, sem agregado}
GROUP BY cols       \Tr{-- 1 linha por grupo}
HAVING COUNT(..) > x   \Tr{-- filtra GRUPOS}
ORDER BY col [ASC | DESC];
\end{Verbatim}
\end{sql}
\textbf{ordem de execução}: FROM \go\ WHERE \go\ GROUP BY \go\ HAVING \go\ SELECT \go\ ORDER BY
\begin{dl}
\item[agregados:] \q{COUNT(*)} conta linhas; \q{COUNT(col)} ignora NULL; \q{COUNT(DISTINCT col)} valores dif.; \q{SUM AVG MIN MAX}
\item[regra:] todo col do SELECT q ñ é agregado vai no GROUP BY; agregado ñ vai no WHERE \go\ HAVING
\item[AVG distrib.:] \q{SUM(s)/SUM(n)} das parciais, nunca AVG de AVG
\item[joins:] \q{INNER JOIN .. ON a = b} só o q casa; \q{NATURAL JOIN} colunas de msm nome; \q{LEFT/RIGHT JOIN} mantém tudo de um lado (NULL no outro) + \q{WHERE x IS NULL} = anti-join
\item[comparação:] \q{= <> < > <= >=}; \q{BETWEEN a AND b} (inclui limites); \q{LIKE 'Fl\%'} (\% qq coisa, \_ 1 char); \q{IS NULL} (nunca \q{= NULL})
\item[subconsulta:] \q{IN (..)}, \q{NOT IN (..)} (NULL na subconsulta \go\ resultado vazio), \q{EXISTS (..)}, \q{> ALL (..)}, \q{= ANY (..)}; no FROM precisa alias: \q{FROM (SELECT ..) P}
\item[conjuntos:] \q{UNION} tira duplicatas; \q{UNION ALL} mantém (+ barato, p/ juntar parciais); \q{INTERSECT}; \q{EXCEPT}
\item[view:] \q{CREATE VIEW v AS SELECT ..}; \q{DROP VIEW v}
\item[DML:] \q{INSERT INTO t (col1, col2) VALUES (v1, v2)}; \q{UPDATE t SET c = v WHERE ..}; \q{DELETE FROM t WHERE ..}
\item[alias:] \q{AS}; Postgres ñ aceita alias do SELECT no HAVING \go\ repete o agregado (\q{HAVING COUNT(*) > 1000})
\end{dl}

\h{SGBDD: JUNTAR CONSULTAS ENTRE NÓS}
\begin{dl}
\item[passo a passo:] \textbf{1.} troca cada tabela global pelos fragmentos (\q{X.T}); \textbf{2.} corta frag. q o WHERE exclui; \textbf{3.} WHERE e só as colunas usadas vão p/ DENTRO de cada nó; \textbf{4.} agrega onde estão os dados; \textbf{5.} junta no nó de destino; \textbf{6.} HAVING / ORDER BY final no destino
\item[view:] \q{CREATE VIEW B.P AS ..} roda no nó B e fica lá; outro nó ler \q{B.P} = tráfego do resultado da view. Nos modelos abaixo: \N{B} = nó remoto, o SELECT final roda em \N{A} (destino)
\end{dl}
\begin{sql}
\begin{Verbatim}
\N{horizontal: une as partes}
CREATE VIEW B.P AS SELECT a, b FROM B.T WHERE c;
SELECT a, b FROM A.T WHERE c
UNION ALL SELECT * FROM B.P;
\N{COUNT/SUM: soma parciais; MIN/MAX: MIN/MAX delas}
CREATE VIEW B.P AS SELECT COUNT(*) AS n,
  SUM(x) AS s, MAX(x) AS mx FROM B.T WHERE c;
SELECT SUM(n), SUM(s), MAX(mx), SUM(s)/SUM(n) AS media
FROM (SELECT COUNT(*) AS n, SUM(x) AS s, MAX(x) AS mx
      FROM A.T WHERE c
      UNION ALL SELECT * FROM B.P) X;
\N{GROUP BY: agrupa em cada no, reagrupa no destino}
CREATE VIEW B.P AS SELECT g, COUNT(*) AS n,
  SUM(x) AS s FROM B.T GROUP BY g;
SELECT g, SUM(n) AS n, SUM(s)/SUM(n) AS media
FROM (SELECT g, COUNT(*) AS n, SUM(x) AS s
      FROM A.T GROUP BY g
      UNION ALL SELECT * FROM B.P) X
GROUP BY g HAVING SUM(n) > 1000 ORDER BY n DESC;
\N{vertical: junta pela chave}
SELECT T1.a, T2.b FROM A.T1 T1
  INNER JOIN B.T2 T2 ON T1.id = T2.id;
\end{Verbatim}
\end{sql}
\begin{dl}
\item[HAVING:] só no destino, dps de somar as parciais. Exceção: frag. pelo atributo do GROUP BY (grupo inteiro num nó só) \go\ HAVING já no nó, manda só os grupos q passam
\item[COUNT(DISTINCT):] ñ soma entre nós se o msm valor pode estar nos 2 \go\ manda os valores distintos e conta no destino
\item[1 nó c/ tudo:] (réplica completa, ex. RS no 4.5) resolve inteiro lá e manda só o resultado
\end{dl}
\end{multicols*}

\newpage
%==================== VERSO: BDD ====================
\fontsize{6.95}{7.7}\selectfont
\begin{multicols*}{3}
\h{BDD: CONCEITOS}
\begin{dl}
\item[BDD:] dados c/ props de BD centralizado (SQL, compartilhado, 1 esquema integrado), mas armazenados/gerenciados em vários locais
\item[SBDD:] nós c/ BD local q juntos agem como 1 BD lógico; consulta local e remota
\item[homog./heterog.:] msm SGBD / SGBDs dif. (até ñ relacional, via gateways)
\item[Date \& Codd:] autonomia local $|$ ñ depender de nó central (ponto único de falha, gargalo) $|$ operação contínua (backup online) $|$ transparência de localização $|$ indep. de fragmentação $|$ indep. de replicação $|$ consultas distribuídas (desempenho indep. de onde submete) $|$ transações distribuídas (ACID) $|$ indep. de HW, SO, rede, SGBD
\item[na prática:] maioria ñ dá indep. total; ACID viola em parte o "ñ nó central" (coordenador)
\item[CAP:] no máx 2 de: Consistência (lê último valor ou erro), Disponibilidade (toda req. responde s/ erro), tolerância a Partição
\item[+:] flexibilidade, autonomia, confiabilidade/disponibilidade (réplicas), desempenho (paralelo, dado perto), escalabilidade
\item[--:] custo de dev, bugs sutis, overhead de msgs, projeto (frag./repl.), conhecimento global difícil (concorrência, deadlock)
\item[replicação:] + paralelismo, + disponib., leitura rápida, update caro
\begin{itemize}
\item \textbf{síncrona}: todas as cópias antes do C (transparente)
\item \textbf{assíncrona}: periódica, pode ficar inconsistente, user sabe (produtos atuais)
\end{itemize}
\item[LAV:] global 1º, fragmentos = visões do global (s/ legado)
\item[GAV:] locais 1º (legado), global = visão integrada (problema de integração)
\end{dl}

\h{FRAGMENTAÇÃO (notação da prova)}
\begin{dl}
\item[horizontal:] \sel\ (tuplas): \q{FragJovens} \la\ \sel$_{age<25}$(TopModels); \q{Rio.TopModels} \la\ \q{FragJovens}
\item[vertical:] \pj\ c/ a \textbf{chave} sempre: \q{SP.TopModels} \la\ \pj$_{mid,age}$(TopModels)
\item[derivada:] R$_i$ \la\ \pj$_{R.*}$(S$_i$ \jn\ R) \ (= R $\ltimes$ S$_i$)
\item[regras:] \textbf{completude} (todo dado em algum frag.), \textbf{reconstrução} (\un\ horiz.; \jn\ pela chave vert.), \textbf{disjunção}
\item[2 passos:] especificação \q{Frag} \la\ expr + atribuição \q{Nodo.Tab} \la\ \q{Frag}
\end{dl}
\sub{Ex. BR} (US igual)
\begin{itemize}
\item \q{BR.Produto} \la\ Produto \ (replicado)
\item \q{BR.Fabricante} \la\ \q{FabBR} \la\ \sel$_{idPais='BR'}$(Fabricante)
\item \q{BR.Capacidade} \la\ \q{CapBR} \la\ \pj$_{C.*}$(FabBR \jn\ Capacidade C)
\item \q{BR.Lote} \la\ \q{LoteBR} \la\ \pj$_{L.*}$(CapBR \jn\ Lote L)
\item \q{BR.Cliente} \la\ \q{CliBR} \la\ \sel$_{idPais='BR'}$(Cliente)
\item \q{BR.Compra} \la\ \pj$_{C.*}$(CliBR \jn\ Compra C) \un\ \pj$_{C.*}$(LoteBR \jn\ Compra C)
\end{itemize}
\sub{Ex. rede social}
\begin{itemize}
\item \q{SC.Pessoa} \la\ \sel$_{uf='SC'}$(Pessoa)
\item \q{SC.Segue} \la\ \pj$_{Segue.*}$(PessoaSC \jn$_{idP=idSeguidor}$ Segue) \un\ idem c/ \q{idSeguido}
\end{itemize}

\h{CONSULTA DISTRIBUÍDA: COMO ESCREVER}
\begin{itemize}
\item resultado vai p/ o nó onde foi \textbf{submetida}
\item cada passo c/ \N{nó} e número (1a, 1b, 2..), usando só tabelas DAQUELE nó ou views recebidas (citar \q{US.X} estando em BR = tráfego)
\item filtra e agrega ANTES de mandar; manda \textbf{chaves}; elimina frag. q o predicado exclui; réplica \go\ usa a + barata
\item \Tr{tráfego} em cada seta: k $\cdot$ (|attr1|+|attr2|). A compensa se tráfego(A) $<$ tráfego(B)
\end{itemize}
\sub{AVG c/ frag. horizontal} (Rio age$<$25, SP age$\ge$25; 20$<$age$<$35, submetida no Rio). Média de médias ERRADA, manda SUM e COUNT
\begin{sql}
\begin{Verbatim}
\N{1a  Rio}
CREATE VIEW Rio.Parc AS
  SELECT SUM(sal) AS s, COUNT(*) AS n
  FROM Rio.TopModels WHERE age > 20;
\N{1b  SP}                       \Tr{-> Rio: |s|+|n|}
CREATE VIEW SP.Parc AS
  SELECT SUM(sal) AS s, COUNT(*) AS n
  FROM SP.TopModels WHERE age < 35;
\N{2  Rio}
SELECT SUM(s)/SUM(n) FROM (SELECT * FROM Rio.Parc
  UNION ALL SELECT * FROM SP.Parc) P;
\end{Verbatim}
\end{sql}
\sub{vertical} (SP = mid,age; Rio = mid,name,sal; no Rio)
\begin{sql}
\begin{Verbatim}
\N{1  SP}                        \Tr{-> Rio: k*|mid|}
CREATE VIEW SP.Parc AS SELECT mid FROM SP.TopModels
  WHERE age > 20 AND age < 35;
\N{2  Rio}
SELECT AVG(sal) FROM Rio.TopModels T
  INNER JOIN SP.Parc P ON T.mid = P.mid;
\end{Verbatim}
\end{sql}
\sub{semi-join} R em N1, S em N2
\begin{enumerate}[label=\textbf{\color{nd}\arabic*}.,leftmargin=3mm]
\item \N{N1}: \q{N1.P} \la\ \pj$_{Key}$(\sel$_{c1}$(N1.R)) \Tr{\go\ N2}
\item \N{N2}: \q{N2.P} \la\ \pj$_{L2}$(\sel$_{c2}$(N2.S) \jn\ N1.P) \Tr{\go\ N1}
\item \N{N1}: \pj$_{L1\cup L2}$(N1.R \jn\ N2.P). Compensa se |chaves| + volta $<$ tabela inteira
\end{enumerate}
\sub{"de fora" c/ frag. derivada} = anti-join local
\begin{sql}
\begin{Verbatim}
SELECT SUM(valor_total) FROM BR.Compra
  WHERE idC NOT IN (SELECT idC FROM BR.Cliente);
\N{ou}
SELECT SUM(Co.valor_total) FROM BR.Cliente C
  RIGHT JOIN BR.Compra Co ON C.idC = Co.idC
  WHERE C.idC IS NULL;
\end{Verbatim}
\end{sql}

\h{DEMONSTRAÇÕES (BR/US, submetidas em BR)}
\begin{dl}
\item[Q1] clientes BR: tudo em BR, s/ tráfego
\item[Q2] fab. BR, cliente $\ne$ BR: anti-join acima, s/ tráfego
\item[Q3] SUM, cliente US e fab. BR:
\end{dl}
\begin{sql}
\begin{Verbatim}
\N{1  BR}  fab. BR q venderam p/ fora  \Tr{-> US: k_BR*|idF|}
CREATE VIEW BR.FbrVendExt AS SELECT DISTINCT Co.idF
  FROM BR.Cliente C RIGHT JOIN BR.Compra Co
  ON C.idC = Co.idC WHERE C.idC IS NULL;
\N{2  US}                       \Tr{-> BR: |valor_total|}
CREATE VIEW US.Resultado AS SELECT SUM(Co.valor_total)
  FROM US.Compra Co NATURAL JOIN BR.FbrVendExt;
\N{3  BR}
SELECT * FROM US.Resultado;
\end{Verbatim}
\end{sql}
\begin{dl}
\item[Q4] = Q3 + \q{GROUP BY F.*}:
\begin{enumerate}[label=\textbf{\color{nd}\arabic*}.,leftmargin=3mm]
\item \N{BR}: \q{BR.FbrVendExt} como Q3 \Tr{(k$_{BR}$|idF|)}
\item \N{US}: \q{CREATE VIEW US.CliCompraBR AS SELECT DISTINCT Co.idC FROM US.Compra Co NATURAL JOIN BR.FbrVendExt} \Tr{(k$_{BR-US}$|idC|)}
\item \N{BR}: \q{SELECT F.*, SUM(Co.valor\_total) FROM US.CliCompraBR NATURAL JOIN BR.Compra Co NATURAL JOIN BR.Lote NATURAL JOIN BR.Capacidade NATURAL JOIN BR.Fabricante F GROUP BY F.*}
\end{enumerate}
alt.: US agrega por idF (\q{SELECT idF, SUM(..) tot .. GROUP BY idF}, \Tr{k(|idF|+|tot|)}) e BR junta c/ \q{BR.Fabricante}
\item[Q5] fab. US, cliente BR, \q{GROUP BY F.*}:
\begin{enumerate}[label=\textbf{\color{nd}\arabic*}.,leftmargin=3mm]
\item \N{BR}: \q{CREATE VIEW BR.CliCompraExt AS SELECT DISTINCT Co.idC FROM BR.Lote L RIGHT JOIN BR.Compra Co ON L.idF=Co.idF WHERE L.idF IS NULL} \Tr{(k$_{BR}$|idC|)}
\item \N{US}: F.*, SUM agrupado juntando \q{US.Compra}, \q{BR.CliCompraExt}, \q{US.Fabricante} \Tr{(k$_{BR-US}$(|F.*|+|valor|))}
\end{enumerate}
alt.: US manda \q{DISTINCT F.*} dos fab. q venderam p/ fora \Tr{(k$_{US}$|F.*|)} e BR agrega. Compensa se k$_{US}$|F.*| $<$ k$_{BR-US}$(|F.*|+|valor|) + k$_{BR}$|idC|
\end{dl}

\h{EXERCÍCIOS RESOLVIDOS}
\sub{Rede social} (E4 Q4 / Simulado Q1)
\begin{itemize}
\item seguidos DO PRÓPRIO estado c/ $>$1M seguidores DE FORA, submetida no msm nó \go\ \ans{só no nó local, s/ tráfego}
\item \q{X.Segue} já tem todo par c/ seguido de X; "de fora" = \q{idSeguidor NOT IN (SELECT idP FROM X.Pessoa)}
\item seguidores de um estado ESPECÍFICO remoto (ex. de MG, submetida em SC): MG filtra e agrega, manda poucos idP:
\end{itemize}
\begin{sql}
\begin{Verbatim}
\N{1  MG}                      \Tr{-> SC: poucos idP}
CREATE VIEW MG.Cand AS SELECT idSeguido, COUNT(*) AS n
  FROM MG.Segue
  WHERE idSeguidor IN (SELECT idP FROM MG.Pessoa)
  AND idSeguido NOT IN (SELECT idP FROM MG.Pessoa)
  GROUP BY idSeguido HAVING COUNT(*) > 1000000;
\N{2  SC}
SELECT P.*, C.n FROM SC.Pessoa P
  INNER JOIN MG.Cand C ON P.idP = C.idSeguido
  ORDER BY C.n DESC;
\end{Verbatim}
\end{sql}
\sub{4.1} LIVRO, LIVRARIA, ESTOQUE
\begin{dl}
\item[a)] \q{LIV1} \la\ \sel$_{Estado='SC'}$(LIVRARIA), \q{LIV2} \la\ \sel$_{Estado\ne'SC'}$(LIVRARIA) (completos e disjuntos)
\item[b)] \q{ESTi} \la\ \pj$_{ESTOQUE.*}$(ESTOQUE \jn\ LIVi)
\item[c)] \q{LVRk} \la\ \sel$_{Assunto='k'}$(LIVRO) p/ cada assunto (+ "outros")
\item[d)] \q{ESTik} \la\ \pj$_{ESTOQUE.*}$(ESTOQUE \jn\ LIVi \jn\ LVRk)
\end{dl}
\sub{4.2} LIVRO por preço: B1 $\le$20, B2 20,01-50, B3 50,01-100, B4 $>$100. LESTE\{B1,B4\} CENTRO\{B1,B2\} OESTE\{B1,B2,B3,B4\}. LIVRARIA por CEP: S1 LESTE, S2 CENTRO, S3 OESTE; ESTOQUE derivado de LIVRARIA
\begin{dl}
\item[a)] \q{preco > 15 AND preco < 55} submetida no LESTE. B1 parte (\q{>15}), B2 inteiro, B3 parte (\q{<55}), \at{B4 sai}. B1 é local; B2 e B3 vêm do \ans{OESTE} (único c/ B3 \go\ 1 nó remoto só; B2 tb poderia vir do CENTRO)
\end{dl}
\begin{sql}
\begin{Verbatim}
\N{1  OESTE}   \Tr{-> LESTE: k*(|Num_livro|+|Estoque_total|)}
CREATE VIEW OESTE.R AS
  SELECT Num_livro, Estoque_total FROM OESTE.B2
  UNION ALL
  SELECT Num_livro, Estoque_total FROM OESTE.B3
  WHERE preco < 55;
\N{2  LESTE}
SELECT Num_livro, Estoque_total FROM LESTE.B1
  WHERE preco > 15
UNION ALL
SELECT * FROM OESTE.R;
\end{Verbatim}
\end{sql}
\begin{dl}
\item[b)] 45\go 55 \go\ o livro sai de B2 e entra em B3. Apaga de TODAS as cópias de B2 (CENTRO, OESTE), insere na única cópia de B3 (OESTE). Tudo na msm transação (2PC). ESTOQUE ñ muda. \at{Valor\_estoque\_total} das livrarias q têm o livro sobe 10$\cdot$Qtd, em cada região
\end{dl}
\begin{sql}
\begin{Verbatim}
\N{1  OESTE}  copia p/ B3 c/ preco novo e apaga de B2
INSERT INTO OESTE.B3
  SELECT Num_livro, Autor_principal, Assunto,
    Estoque_total, 55.00
  FROM OESTE.B2 WHERE Num_livro = 1234;
DELETE FROM OESTE.B2 WHERE Num_livro = 1234;
\N{2  CENTRO}  a outra copia de B2
DELETE FROM CENTRO.B2 WHERE Num_livro = 1234;
\N{3  LESTE, CENTRO e OESTE}  LIVRARIA/ESTOQUE locais
UPDATE LIVRARIA L SET Valor_estoque_total =
  Valor_estoque_total + 10 * (SELECT E.Qtd
    FROM ESTOQUE E WHERE E.Num_livro = 1234
    AND E.Num_livraria = L.Num_livraria)
WHERE L.Num_livraria IN (SELECT Num_livraria
  FROM ESTOQUE WHERE Num_livro = 1234);
\end{Verbatim}
\end{sql}
\begin{dl}
\item[c)] qualquer coisa q precise de S2 (ou do ESTOQUE derivado dele), q só existe no CENTRO. Ex. no OESTE: \q{SELECT Num\_livraria, Cidade FROM LIVRARIA WHERE CEP BETWEEN 35001 AND 70000}
\end{dl}
\sub{4.3} Funcionario(codF, nome, nroFilial) no nodo1; Filial Floripa(20) nodo2, Blumenau(10) nodo3; submetida no nodo2
\begin{dl}
\item[a)] \q{N1.Funcionario} \la\ Funcionario; \q{N2.Filial} \la\ \sel$_{cidade='Florianopolis'}$(Filial); \q{N3.Filial} \la\ \sel$_{cidade='Blumenau'}$(Filial)
\item[b) Alt1:] filtra no nodo1, manda só o resultado
\end{dl}
\begin{sql}
\begin{Verbatim}
\N{1a  nodo2}                  \Tr{-> nodo1: 20*|nroFilial|}
CREATE VIEW N2.FilFlp AS
  SELECT nroFilial FROM N2.Filial;
\N{1b  nodo3}                  \Tr{-> nodo1: 10*|nroFilial|}
CREATE VIEW N3.FilBlu AS
  SELECT nroFilial FROM N3.Filial;
\N{2  nodo1}        \Tr{-> nodo2: k*(|codF|+|nome|+|nroF|)}
CREATE VIEW N1.Res AS
  SELECT codF, nome, nroFilial FROM N1.Funcionario
  WHERE nroFilial IN (SELECT nroFilial FROM N2.FilFlp
    UNION SELECT nroFilial FROM N3.FilBlu);
\N{3  nodo2}
SELECT * FROM N1.Res;
\end{Verbatim}
\end{sql}
\begin{dl}
\item[b) Alt2:] traz todos os func. p/ o nodo2 e resolve lá
\end{dl}
\begin{sql}
\begin{Verbatim}
\N{1a  nodo1}       \Tr{-> nodo2: N*(|codF|+|nome|+|nroF|)}
CREATE VIEW N1.Todos AS
  SELECT codF, nome, nroFilial FROM N1.Funcionario;
\N{1b  nodo3}                  \Tr{-> nodo2: 10*|nroFilial|}
CREATE VIEW N3.FilBlu AS
  SELECT nroFilial FROM N3.Filial;
\N{2  nodo2}
SELECT F.codF, F.nome, F.nroFilial FROM N1.Todos F
  WHERE F.nroFilial IN (SELECT nroFilial FROM N2.Filial
    UNION SELECT nroFilial FROM N3.FilBlu);
\end{Verbatim}
\end{sql}
\begin{dl}
\item[c)] Alt1 \Tr{30$\cdot$|nroF| + k$\cdot$(|codF|+|nome|+|nroF|)}; Alt2 \Tr{10$\cdot$|nroF| + N$\cdot$(|codF|+|nome|+|nroF|)}; k = func. das 30 filiais, N = todos
\item[d)] \ans{Alt1 ganha} qdo k $\ll$ N (normal, piora p/ Alt2 c/ a empresa crescendo p/ outras cidades); empata só se quase todo func. é de Floripa/Blumenau
\end{dl}
\sub{4.5} Aluno/Cursa/Disciplina; SC frag., RS c/ tudo replicado; submetida em SC
\begin{dl}
\item[a)] \q{SC.Aluno} \la\ \sel$_{uf='SC'}$(Aluno); \q{SC.Cursa} \la\ \pj$_{C.*}$(AlunoSC \jn\ Cursa C); \q{SC.Disciplina} \la\ Disciplina; \q{RS.*} \la\ tabelas inteiras
\end{dl}
\begin{sql}
\begin{Verbatim}
\N{1  RS}            \Tr{-> SC: k*(|idD|+|nAlunos|+|mNotas|)}
CREATE VIEW RS.Res AS
  SELECT idD, COUNT(DISTINCT idA) AS nAlunos,
  AVG(nota) AS mNotas FROM RS.Cursa GROUP BY idD
  HAVING COUNT(DISTINCT idA) > 1000 AND AVG(nota) < 7;
\N{2  SC}  D.* local, Disciplina replicada
SELECT D.*, R.nAlunos, R.mNotas FROM SC.Disciplina D
  NATURAL JOIN RS.Res R
  ORDER BY mNotas ASC, nAlunos DESC;
\end{Verbatim}
\end{sql}
\at{ñ usar SC.Cursa + RS}: contaria aluno de SC 2x

\h{2PC (two-phase commit) = ATOMICIDADE}
\begin{dl}
\item[fase 1:] coord. grava \q{<PREPARE T>} e pede p/ todos se prepararem. Pronto \go\ grava \q{<READY T>}, avisa coord. e \at{ñ pode mais abortar sozinho}. Ñ consegue \go\ responde ñ (ou ñ responde = timeout)
\item[fase 2:] TODOS READY \go\ COMMIT p/ todos. Algum ñ ou timeout \go\ ABORT p/ todos
\item[coord. cai:] dps dos READY \go\ participantes \textbf{bloqueados} (3PC reduz). Nó READY q volta pergunta ao coord.
\end{dl}
\sub{RECEITA: nó X falha $>$ timeout da fase 1}
\aviso{\at{colunas = logs independentes, linhas ñ são tempo global}}
\begin{enumerate}[label=\textbf{\color{hd}\arabic*.},leftmargin=3mm]
\item X participa de T? (tem \q{<START T>}) ñ \go\ T segue normal
\item X já tinha \q{<READY T>} (ou \q{<PREPARE T>} se é o coord.) antes de cair? \textbf{sim} \go\ voto dele conta:
\begin{itemize}
\item outros tb READY \go\ \ans{COMMIT}
\item outros ainda ñ: "outros continuam normalmente" \go\ \ans{COMMIT}; enunciado ñ diz \go\ \ans{"COMMIT ou ABORT, depende dos outros"}
\end{itemize}
\item X caiu \textbf{s/} READY/PREPARE de T \go\ \ans{ABORT}
\end{enumerate}
\sub{Respostas}
\begin{dl}
\item[E4:] BR cai (tinha READY T2): T2 \go\ \ans{COMMIT}\\
UK cai (s/ READY T1): T1 \go\ \ans{ABORT}\\
variante (\q{RL2(B) WL2(C) R2(B) UL2(B)..}) \go\ \ans{conserv. e estrito}
\item[Simulado:] Q1 rede social \go\ \ans{só SC, s/ tráfego}\\
Q2 coord. BR cai c/ PREPARE T2, US/UK s/ READY T2 \go\ \ans{COMMIT ou ABORT, depende} (prov.)\\
Q3 variante \go\ \ans{conserv. e estrito}\\
Q4 US cai s/ READY T1 \go\ \ans{ABORT}
\item[4.4:] SP cai: T1 \go\ \ans{ABORT}, T2 \go\ \ans{COMMIT}\\
SC cai, RJ coordena: T1 \go\ \ans{ABORT}, T2 \go\ \ans{COMMIT} (prov.)\\
RJ cai, outros normais: T1 \go\ \ans{COMMIT}, T2 \go\ \ans{COMMIT}\\
variante \go\ \ans{ñ conserv. e ñ estrito}
\end{dl}
\begin{itemize}
\item 4.4 logs: RJ \q{RL1(A) R1(A) WL1(B) UL1(A) W1(B) UL1(B) <READY T1>}; SP \q{RL2(C) R2(C) WL2(C) W2(C) UL2(C) <READY T2>}, s/ READY T1; SC coord \q{RL2(D) R2(D) UL2(D) <PREPARE T2> WL1(D) W1(D)}, s/ PREPARE T1
\item variante 4.4: WL1(B) dps de R1(A) e upgrade RL2\go WL2(C) \go\ ñ conserv.; UL de WL antes do C \go\ ñ estrito
\end{itemize}

\h{CONCORRÊNCIA / DEADLOCK EM BDD}
\begin{dl}
\item[otimista:] execução livre, monitora conflitos (grafos ou relógios lógicos); conflito \go\ desfaz a + nova
\item[lock:] troca de msgs p/ pedir bloqueio local; deadlock + difícil de detectar
\item[cópia primária:] coordena os locks do item replicado
\item[cópia secundária:] atualizada síncrona c/ a primária, assume se ela cai
\item[site primário:] (coord. único) todas as primárias no msm site; simples, mas ponto único de falha e gargalo
\item[coord. múltiplos:] primárias em sites dif.
\item[eleição:] escolhe novo coord. se o atual cai
\item[protocolo desviado:] RL só local, WL em TODAS as réplicas
\item[votação:] maioria das réplicas
\item[marcador de tempo:] centralizado (relógio/contador) ou distribuído (relógios lógicos)
\item[deadlock distrib.:] ciclo pode existir só no grafo GLOBAL (T1\go T2 no site A, T2\go T1 no B)
\begin{itemize}
\item \textbf{centralizada}: grafos locais vão p/ 1 nó q junta e procura ciclo
\item \textbf{hierárquica}: cada nó manda p/ o pai
\item \textbf{timeout}: espera demais \go\ aborta e recomeça
\end{itemize}
\end{dl}
\end{multicols*}
\end{document}
